%% ============================================================
%% Lineamientos del Proyecto Final
%% Analítica para los negocios — 06327-ECO
%% Universidad Icesi | Periodo 2026-02
%% ============================================================
\documentclass[11pt,letterpaper]{article}

\usepackage[utf8]{inputenc}
\usepackage[T1]{fontenc}
\usepackage[top=2.5cm, bottom=2.5cm, left=2.8cm, right=2.8cm]{geometry}
\usepackage{parskip}
\usepackage{booktabs}
\usepackage{tabularx}
\usepackage{array}
\usepackage{enumitem}
\usepackage{fancyhdr}
\usepackage{titlesec}
\usepackage{hyperref}
\usepackage{lastpage}data:image/png;base64,iVBORw0KGgoAAAANSUhEUgAAABIAAAASCAYAAABWzo5XAAAAbElEQVR4Xs2RQQrAMAgEfZgf7W9LAguybljJpR3wEse5JOL3ZObDb4x1loDhHbBOFU6i2Ddnw2KNiXcdAXygJlwE8OFVBHDgKrLgSInN4WMe9iXiqIVsTMjH7z/GhNTEibOxQswcYIWYOR/zAjBJfiXh3jZ6AAAAAElFTkSuQmCC

\hypersetup{
  colorlinks=true,
  linkcolor=black,
  urlcolor=blue,
  citecolor=black
}
%% dejo comentado lo que voy cambiando 
%% para que lo veas y me comentes
%% si te parece bien o cambiarías algo

%% ─── Secciones ───────────────────────────────────────────────
\titleformat{\section}{\large\bfseries}{}{0em}{}[\vspace{2pt}\hrule\vspace{4pt}]
\titleformat{\subsection}{\normalsize\bfseries}{}{0em}{}
\titleformat{\subsubsection}{\normalsize\itshape\bfseries}{}{0em}{}
\titlespacing{\section}{0pt}{18pt}{6pt}
\titlespacing{\subsection}{0pt}{12pt}{4pt}
\titlespacing{\subsubsection}{0pt}{8pt}{2pt}

%% ─── Encabezado y pie ────────────────────────────────────────
\pagestyle{fancy}
\fancyhf{}
\fancyhead[L]{\footnotesize\itshape Analítica para los Negocios $|$ Proyecto Final $|$ U. Icesi 2026-02}
\fancyhead[R]{\footnotesize Pág.\ \thepage\ de \pageref{LastPage}}
\fancyfoot[C]{\footnotesize\itshape Lineamientos del Proyecto Final --- uso exclusivo del curso}
\renewcommand{\headrulewidth}{0.4pt}
\renewcommand{\footrulewidth}{0.4pt}

%% ─── Columnas tabla ──────────────────────────────────────────
\newcolumntype{L}[1]{>{\raggedright\arraybackslash}p{#1}}
\newcolumntype{C}[1]{>{\centering\arraybackslash}p{#1}}
\setlength{\tabcolsep}{6pt}
\renewcommand{\arraystretch}{1.35}

\setlist[itemize]{label=\textbullet, leftmargin=1.5em, itemsep=2pt, topsep=4pt}
\setlist[enumerate]{leftmargin=1.8em, itemsep=3pt, topsep=4pt}

%% ============================================================
\begin{document}

%% ─── PORTADA ─────────────────────────────────────────────────
\thispagestyle{empty}
\vspace*{2.5cm}

\begin{center}
  {\large\textbf{ANALÍTICA PARA LOS NEGOCIOS}}\\[5pt]
  {\normalsize 06327-ECO\ $|$\ Universidad Icesi\ $|$\ Periodo 2026-02}
\end{center}

\vspace{3cm}
\begin{center}
  {\LARGE\bfseries Lineamientos del Proyecto Final}\\[10pt]
  {\normalsize\itshape Guía de entregas, requisitos mínimos, formato y rúbrica de evaluación}
\end{center}

\vspace{3.5cm}
\begin{center}
\begin{tabular}{@{}ll@{}}
  \textbf{Profesor}           & Eduard Martínez González \\[3pt]
  \textbf{Grupos}             & 3 estudiantes por grupo \\[3pt]
  \textbf{Peso evaluativo}    & 25\% de la nota final \\[3pt]
  \textbf{Número de entregas} & 3 entregas progresivas \\
\end{tabular}
\end{center}
%% Se cambian los porcentajes; recuerda integrar
%% con los demas: E1 (0%), E2 (5%) y E3 (20%)
\clearpage
\tableofcontents
\clearpage

%% ============================================================
\section{1. Descripción general del proyecto}

El proyecto final es el eje integrador del curso; su propósito es que el grupo enfrente una situación de negocio real, tome decisiones analíticas justificadas, ejecute el análisis usando las herramientas del curso y comunique resultados de manera clara y accionable.

Se entregará a los estudiantes un \textbf{caso general con una base de datos amplia} (sección 2). Cada grupo tendrá la potestad de \textbf{definir su propia línea de análisis} y deberá mantenerla a lo largo de las tres entregas del semestre. Los grupos pueden optar por una línea de \textbf{clasificación supervisada, segmentación no supervisada o regresión}; en cualquiera de los tres casos, deberán justificar por qué eligieron ese enfoque.

%% Cada grupo \textbf{define un problema} relacionado a los datos proporcionados y trabaja alrededor de este a lo largo de las tres entregas del semestre. pueden hacer \textbf{aprendizaje supervisado} y \textbf{no supervisado}.

El foco del proyecto \textbf{no es escribir código perfecto}, es interpretar resultados, tomar decisiones justificadas y comunicar de manera que alguien sin conocimientos técnicos pueda entender y actuar.

\subsection{Conformación de grupos}

El proyecto se desarrolla en grupos de exactamente \textbf{3 estudiantes}. Cada grupo debe conformarse antes de la fecha de la Entrega~1 y definir el problema sobre el que van a trabajar: \textbf{no se permiten cambios después de la Entrega~1.}

Cabe destacar que, a partir de la segunda entrega, se priorizará la comunicación de los resultados. Por lo tanto, \textbf{no se requerirá un documento escrito, sino que las entregas subsiguientes se evaluarán exclusivamente mediante el soporte visual (diapositivas) y la sustentación oral.}

%% ============================================================
%% otra sección para que revises y me dices que opinas

%% ============================================================
%% otra sección para que revises y me dices que opinas
%% añadi las especificaciones sobre el tipo de entregas y la
%% prioridad al momento de evaluar (Bolaños)
%% ============================================================
\subsection{Importante}

Es de carácter obligatorio enviar los entregables en las fechas establecidas. El incumplimiento de esta directriz sin una excusa válida resultará en una \textbf{calificación de 0.0 para la actividad correspondiente, además de inhabilitar las siguientes entregas.}

Por su parte, la presentación final equivale a un examen parcial y se rige por la sección XIII del \href{https://www.icesi.edu.co/wp-content/uploads/2024/08/Reglamento-estudiantes-pregrado-LDDN-marzo-2026.pdf}{ \textbf{Libro de derechos, deberes y normativa de los estudiantes de pregrado}}. Por lo tanto, su fecha se encuentra definida desde el inicio del periodo académico; si no se puede asistir, deberá tramitar la solicitud de un examen supletorio siguiendo el conducto regular.
%% ============================================================
%% aqui cierra
%% se mejoro la redacción, ademas de incluir la penalización
%% de no entregar las entregas anteriores
%% ============================================================
\subsection{Regla de autoría y uso de IA}

El proyecto debe ser resultado del trabajo del grupo. El uso de IA está permitido como apoyo, con las siguientes restricciones: no se puede entregar como propio un análisis o texto generado por IA sin citarlo ni explicarlo; si usan IA para apoyar el análisis o la redacción, deben indicarlo explícitamente señalando qué herramienta usaron y para qué. \textbf{Regla de oro:} si no pueden explicar y reproducir, no cuenta; en la sustentación se podrá preguntar sobre cualquier decisión del informe; no poder explicarla equivale a no haberla hecho.

En términos de la escala institucional de niveles de uso de la IAG (ver el syllabus del curso): la Entrega~1 corresponde al \textbf{nivel 2 --- Planificación con IAG} (la IA puede usarse para explorar ideas y perspectivas, pero la pregunta entregada debe ser formulada y redactada por el propio grupo, sin texto generado por IA), mientras que las Entregas~2 y~3 corresponden al \textbf{nivel 3 --- Colaboración con IAG} (pueden apoyarse en la IA para desarrollar los entregables, llevando registro de las interacciones y los \textit{prompts}, y demostrando en la sustentación que comprenden y pueden modificar todo lo producido).

%% ============================================================
%% dejo comentada la sección porque toca revisar qué hacer
%% en ella, si borrar todo o reemplazar
%% ============================================================
\clearpage
\section{2. Presentación del caso disponible}

\subsection{El caso Cóndor}
\textbf{La empresa}. Cóndor es una fintech andina con sede en Colombia y operación incipiente en Perú y Ecuador. Opera una plataforma de pagos y crédito de dos lados: del lado del consumidor, una billetera digital que permite guardar dinero, pagar con QR y tarjeta, transferir, recargar y solicitar microcréditos; del lado del comercio, datáfonos y códigos QR para aceptar pagos, junto con adelantos de capital de trabajo calculados a partir del flujo de ventas de cada negocio.

\textbf{Crecimiento}. En cinco años de expansión acelerada, Cóndor pasó de ser una empresa emergente a procesar millones de transacciones y a conectar a consumidores y comercios en tres países. Ese crecimiento, no obstante, ocurrió de forma fragmentada: cada área del negocio —crédito, fraude, marketing, operaciones, finanzas y servicio al cliente— registró su operación por separado, de modo que la compañía acumuló un volumen considerable de datos sin una mirada integral que los articulara.

\textbf{Dolores del negocio}. La Gerencia General enfrenta hoy una preocupación de fondo: la empresa gana usuarios, pero pierde rentabilidad por varios frentes sin lograr precisar dónde ni por qué. Las pérdidas por fraude crecen más rápido que el volumen de operaciones; la cartera de microcrédito registra una mora creciente; el costo de adquisición de clientes aumenta mientras muchos usuarios se tornan inactivos con rapidez; numerosos comercios se activan y luego dejan de transaccionar sin previo aviso; y el área financiera no consigue proyectar con confianza los ingresos para la próxima ronda de inversión.

\textbf{La tarea del grupo}. Frente a este panorama, cada grupo asume el papel de un equipo externo de analistas contratado para elaborar un diagnóstico integral y convertir los datos dispersos en decisiones. El proyecto se desarrolla en equipos de tres integrantes y se estructura en tres entregas progresivas que, en conjunto, representan el 25\% de la nota del curso. Cada grupo selecciona un frente —uno de los departamentos de la empresa—, formula su propia pregunta de negocio y la resuelve mediante una de las tres rutas analíticas del curso —clasificación, regresión o segmentación—, ruta que deberá justificar y mantener a lo largo de todo el proyecto. El propósito no es escribir código perfecto, sino plantear una buena pregunta, elegir la tarea analítica adecuada, interpretar los resultados y comunicarlos de manera que alguien sin conocimientos técnicos pueda comprender y actuar.

El contexto completo de la empresa —el detalle de cada departamento, la estructura de su base de datos y las reglas de análisis— se presenta en el \href{https://eduard-martinez.github.io/teaching/ba/final_project/contexto_caso_condor.pdf} {\textbf{documento de contexto del caso}}.

%% ============================================================
\clearpage
\section{3. Entrega 1 --- formulación de la pregunta de investigación (0\%)}

La Entrega~1 es un documento corto donde el grupo declara el problema y línea de análisis que va a tratar, formula la pregunta de negocio en sus propias palabras y describe el equipo con roles asignados. \textbf{Si bien esta fase inicial no es calificable, constituye un prerrequisito obligatorio para avanzar y presentar las siguientes entregas del proyecto}. El formato de entrega es PDF con un \textbf{máximo de 3 páginas} sin contar portada. La entrega se realiza en la \textbf{semana~6} (del 31 de agosto al 5 de septiembre de 2026), a través de la plataforma Intu.



\subsection{Contenido mínimo requerido}

El documento debe incluir las siguientes secciones; no se aceptan entregas que omitan alguna.

\begin{enumerate}
  \item \textbf{Portada.} Nombre del grupo, nombres completos de los integrantes, formulación del problema y fecha.
  %%\item \textbf{problema y justificación.} Indicar el contexto del problema
 \item \textbf{Linea de análisis.} Justificar en 2-3 oraciones la elección de este enfoque metodológico para el proyecto.
  \item \textbf{Pregunta de negocio en sus palabras.} Reformulación propia, clara, accionable y traducible a una tarea analítica.
  \item \textbf{Traducción a tarea analítica.} Indicar el tipo de tarea, qué variable se predice o agrupa y qué salida esperan obtener.
  \item \textbf{Roles del equipo.} Asignar a cada integrante al menos un rol (Data Engineer, Data Analyst, Data Scientist, Data Translator) con responsabilidades concretas.
  \item \textbf{Variables de interés (preliminar).} Listar al menos 5 variables del dataset que consideran relevantes, con justificación breve.
\end{enumerate}

Esta entrega \textbf{no requiere análisis de datos}. Su propósito es que el grupo plantee la pregunta, analice el problema y planifique el trabajo.

\subsection{Rúbrica --- Entrega 1}
%% ============================================================
%% toca revisar toda la rúbrica
%% verificar si se puede resumir mas la rubrica para que quede
%% en una misma pagina
%% Como la Entrega 1 no tiene calificacion
%% ¿Vale la pena dejar los criterios?
%% ============================================================
\begin{tabular}{L{3.2cm} L{3.4cm} L{3.4cm} L{4.2cm}}
  \toprule
  \textbf{Criterio} & \textbf{Completo} & \textbf{Parcial} & \textbf{Insuficiente} \\
  \midrule
  Línea de análisis
    & Identifica con claridad el enfoque elegido (supervisado, no supervisado o regresión) y justifica por qué es adecuado y viable para su pregunta.
    & Plantea el enfoque pero no justifica por qué es el adecuado para su pregunta.
    & Ausente o copiada del enunciado. \\
  \addlinespace
  Pregunta de negocio
    & Clara, accionable, correctamente planteada.
    & Formulada pero imprecisa o mal clasificada.
    & Ausente o copiada del enunciado. \\
  \addlinespace
  Roles con responsabilidades
    & Cada rol: integrante asignado y descripción concreta.
    & Roles asignados pero sin descripción real.
    & Ausente o roles copiados sin personalización. \\
  \addlinespace
  Variables justificadas
    & $\geq$5 variables con justificación sustantiva.
    & $<$5 variables o justificaciones superficiales.
    & Ausente o solo copia la lista del dataset. \\
  \addlinespace
  Formato y extensión
    & Cumple máximo de páginas y secciones completas.
    & Ligero exceso o secciones incompletas.
    & Supera extensión en $>$50\% o faltan secciones. \\
  \bottomrule
\end{tabular}

%% ============================================================
\clearpage
\section{4. Entrega 2 --- Base de datos depurada y EDA (5\%)}

La Entrega~2 evidencia que el grupo comprendió, limpió y procesó adecuadamente los datos, y realizó un análisis exploratorio que caracteriza patrones relevantes antes de la fase de modelado. Tiene un peso de \textbf{5\% de la nota final del curso}.

\textbf{Formato de entrega.} La evaluación se basa en la \textbf{sustentación oral} apoyada en una presentación de diapositivas (máximo recomendado: \textbf{8 a 10 diapositivas}). Además, el grupo debe adjuntar el \textbf{script de R} (\texttt{.R}, \texttt{.Rmd} o \texttt{.qmd}) usado para la limpieza y el análisis, como soporte técnico verificable. La sustentación se realiza en la \textbf{semana~9} (del 21 al 26 de septiembre de 2026).

%%========================================
%% Preguntar si se deja un limite de slides
%%========================================

\subsection{Contenido mínimo requerido}

\subsubsection{Debe estar realizado (en el script)}

\begin{enumerate}
  \item Diccionario de variables: nombre, tipo y descripción breve de cada una.
  \item Tabla de valores faltantes por variable (número y \%).
  \item Tabla completa de estadísticas descriptivas (media, mediana, SD, mín, máx) de las variables numéricas.
  \item Exploración de la distribución de las variables clave (histogramas o boxplots).
  \item Tratamiento de faltantes y atípicos, con las decisiones tomadas y su justificación.
\end{enumerate}

\subsubsection*{Debe presentarse y sustentarse (en diapositivas)}

\begin{itemize}
  \item Diapositiva de título: mismo grupo y caso de la Entrega~1.
  \item Contexto de los datos: número de observaciones y variables, unidad de análisis y variables clave.
  \item Una tabla-resumen sintética con los descriptivos y faltantes \textit{relevantes} (no la tabla completa).
  \item Mínimo 2 visualizaciones de distribución, priorizando las más informativas.
  \item Un gráfico de la variable objetivo: distribución con \% por clase (clasificación), histograma o densidad de la variable continua (regresión), o distribución de las variables de agrupación (segmentación).
  \item Un gráfico que muestre la relación entre una variable explicativa y el objetivo: comparación por clase o segmento (clasificación/segmentación) o diagrama de dispersión frente a la variable objetivo (regresión), que motive el modelado.
  \item Mínimo 3 hallazgos en lenguaje de negocio y las principales limitaciones identificadas.
\end{itemize}

Todo gráfico presentado debe estar etiquetado y su interpretación será exigida durante la sustentación; un gráfico mostrado sin interpretación oral no tiene validez.

\subsection{Rúbrica --- Entrega 2}

\begin{tabular}{L{2.6cm} L{3cm} L{2.6cm} L{2.6cm} L{2.6cm}}
  \toprule
  \textbf{Criterio} & \textbf{Excelente (100\%)} & \textbf{Bueno (80\%)} & \textbf{Suficiente (60\%)} & \textbf{Insuf.\ (0\%)} \\
  \midrule
  Calidad y limpieza de datos
    & Tratamiento de faltantes y atípicos completo; decisiones justificadas en la exposición.
    & Limpieza completa pero alguna decisión sin justificar.
    & Limpieza incompleta o poco clara.
    & Ausente o incorrecta. \\
  \addlinespace
  EDA y hallazgos de negocio
    & Descriptivos correctos y $\geq$3 hallazgos claros en lenguaje de negocio.
    & EDA presente; hallazgos algo superficiales.
    & EDA mínimo; hallazgos genéricos.
    & EDA ausente o sin análisis. \\
  \addlinespace
  Selección y soporte visual
    & Presenta solo lo esencial; la elección de gráficos evidencia criterio analítico y no satura las diapositivas.
    & Cumple los mínimos; alguna visual accesoria o interpretación débil.
    & No cumple 1--2 mínimos o satura las diapositivas.
    & Cumple menos de la mitad. \\
  \addlinespace
  Código R (reproducibilidad)
    & Script adjunto, funcional y comentado; reproduce lo mostrado.
    & Funciona pero poco comentado.
    & Parcialmente funcional.
    & No se entregó o no corre. \\
  \addlinespace
  Comunicación oral y diapositivas
    & Exposición clara y estructurada; diapositivas limpias; buen manejo del tiempo.
    & Ordenada con transiciones o tiempos débiles.
    & Difícil de seguir o diapositivas saturadas.
    & Desorganizada o incoherente. \\
  \bottomrule
\end{tabular}

%% ============================================================
%% Posiblemente esta sección este sujeta a cambios, leela y
%% evalua si los criterios te parecen correctos.
%% ============================================================
\clearpage
\section{5. Entrega 3 --- Análisis final y sustentación (20\%)}

La Entrega~3 es la culminación del proyecto: el grupo presenta el análisis completo (clasificación, segmentación o regresión, resultados y recomendaciones) y lo sustenta oralmente ante el profesor. No se entrega documento escrito. El peso total es \textbf{20\% de la nota del curso}, distribuido así: \textbf{sustentación oral 15\%}, \textbf{diapositivas 4\%} y \textbf{script de R 1\%}.

\textbf{Formato de entrega.} Presentación de diapositivas (máximo recomendado: \textbf{10 a 12 diapositivas}) sustentada oralmente, más el \textbf{script de R} (\texttt{.R}, \texttt{.Rmd} o \texttt{.qmd}) que reproduzca todos los resultados mostrados. La sustentación se realizará en la \textbf{semana~16} (del 9 al 14 de noviembre de 2026); el profesor puede seleccionar de forma aleatoria a los integrantes que sustentan.

\subsection{Contenido mínimo}

Igual que en la Entrega~2, se distingue lo que debe estar \textbf{realizado} (verificable en el script) de lo que debe \textbf{presentarse y sustentarse}. Las secciones marcadas con (C) aplican solo a clasificación, las marcadas con (S) solo a segmentación y las marcadas con (R) solo a regresión.

\subsubsection*{Debe estar realizado (en el script)}

\begin{itemize}
  \item Preparación final de los datos: limpieza, transformaciones o nuevas variables, sin fugas (\textit{leakage}).
  \item (C) División train/test, entrenamiento de al menos dos modelos y cálculo de métricas.
  \item (R) División train/test, entrenamiento de al menos dos modelos y cálculo de métricas de error.
  \item (S) Estandarización de variables, prueba de al menos dos valores de $k$ y construcción de perfiles.
  \item Generación de todos los resultados, tablas y gráficos que respaldan la presentación.
\end{itemize}

\subsubsection*{Debe presentarse y sustentarse (en diapositivas)}

\begin{itemize}
  \item Diapositiva de título y pregunta de negocio: qué problema se resuelve y qué decisión se toma.
  \item Resumen del EDA: solo los hallazgos que orientan el modelado.
  \item (C) Tabla de métricas (accuracy, precisión, recall, F1) comparada con un \textit{baseline}, y matriz de confusión interpretada en lenguaje de negocio.
  \item (C) Un gráfico de las variables más importantes (``señales''); curva ROC con AUC recomendada.
  \item (R) Tabla de métricas de error (RMSE, MAE y $R^2$) comparada con un \textit{baseline} (por ejemplo, predecir la media), interpretada en lenguaje de negocio.
  \item (R) Un gráfico de valores predichos frente a observados (o de residuos) y un gráfico de las variables más importantes.
  \item (S) Tabla de perfiles: nombre, tamaño (N y \%) y estadísticos clave por segmento.
  \item (S) Un gráfico de \textit{elbow} o \textit{silhouette} indicando el $k$ elegido y mínimo una visualización de los segmentos (PCA/\textit{scatter}, barras o radar).
  \item Interpretación y recomendación: qué acción concreta se recomienda y a qué grupo va dirigida.
  \item Limitaciones y conclusiones: qué no pudo hacerse y qué se mejoraría con más datos.
\end{itemize}

Todo gráfico o tabla mostrado debe interpretarse durante la sustentación; un elemento sin interpretación oral no tiene validez.

\subsection{Rúbrica --- Sustentación oral (15\%)}

\begin{tabular}{C{0.7cm} L{2.7cm} L{2.6cm} L{2.6cm} L{2.6cm} L{2.4cm}}
  \toprule
  \textbf{\%} & \textbf{Criterio} & \textbf{Excelente (100\%)} & \textbf{Bueno (75\%)} & \textbf{Suficiente (50\%)} & \textbf{Insuf.\ (0\%)} \\
  \midrule
  35 & Dominio del análisis
    & Explica con claridad qué hizo, por qué y qué significan los resultados; responde con precisión.
    & Explica bien lo principal, con dudas en detalles.
    & Explica parcialmente; dificultades en varias preguntas.
    & No puede explicar el análisis. \\
  \addlinespace
  25 & Interpretación y recomendación
    & Conecta resultados con una decisión concreta; recomendación justificada.
    & Recomendación presente pero poco justificada.
    & Recomendación muy general.
    & Sin recomendación o desconectada. \\
  \addlinespace
  20 & Respuesta a preguntas
    & Responde con seguridad y coherencia.
    & Responde bien la mayoría.
    & Dificultades ante preguntas de profundidad.
    & No responde o respuestas incoherentes. \\
  \addlinespace
  20 & Claridad y manejo del tiempo
    & Exposición lógica, bien estructurada y dentro del tiempo.
    & Ordenada con  algunas transiciones débiles.
    & Difícil de seguir el hilo.
    & Desorganizada o fuera de tiempo. \\
  \bottomrule
\end{tabular}

\subsection{Rúbrica --- Diapositivas (4\%)}

\begin{tabular}{C{0.7cm} L{2.7cm} L{2.6cm} L{2.6cm} L{2.6cm} L{2.4cm}}
  \toprule
  \textbf{\%} & \textbf{Criterio} & \textbf{Excelente (100\%)} & \textbf{Bueno (75\%)} & \textbf{Suficiente (50\%)} & \textbf{Insuf.\ (0\%)} \\
  \midrule
  40 & Completitud del análisis
    & Pregunta, EDA, modelo/segmentación y resultados presentes y coherentes.
    & Presente con errores menores.
    & Incompleto o con errores conceptuales.
    & Análisis ausente o incorrecto. \\
  \addlinespace
  35 & Selección y calidad visual
    & Muestra solo lo esencial; gráficos legibles, etiquetados y bien elegidos.
    & Cumple mínimos; alguna visual accesoria.
    & Diapositivas saturadas o faltan elementos clave.
    & Soporte visual deficiente. \\
  \addlinespace
  25 & Interpretación de negocio
    & Resultados y recomendación expresados para un decisor no técnico.
    & Interpretación parcialmente técnica.
    & Interpretación superficial.
    & Sin interpretación de negocio. \\
  \bottomrule
\end{tabular}

\subsection{Rúbrica --- Script de R (1\%)}

\begin{tabular}{C{0.7cm} L{2.7cm} L{2.6cm} L{2.6cm} L{2.6cm} L{2.4cm}}
  \toprule
  \textbf{\%} & \textbf{Criterio} & \textbf{Excelente (100\%)} & \textbf{Bueno (75\%)} & \textbf{Suficiente (50\%)} & \textbf{Insuf.\ (0\%)} \\
  \midrule
  50 & Funciona y reproduce
    & Corre sin errores y reproduce todos los resultados mostrados.
    & Corre con ajustes menores.
    & Reproduce solo parte de los resultados.
    & No corre o no se entregó. \\
  \addlinespace
  50 & Orden y comentarios
    & Código limpio, comentado y fácil de seguir.
    & Comentado parcialmente.
    & Poco comentado o desordenado.
    & Sin comentarios ni estructura. \\
  \bottomrule
\end{tabular}
%% ============================================================
%% Verificar si es coherente este apartado; aun me genera dudas
%% la distribucion de los porcentaje y la rubrica aunque me
%% convence, me gustaria una segunda opinion.
%% ============================================================
\clearpage
\section{Resumen de pesos y fechas}

El proyecto final representa el 25\% de la nota del curso y se distribuye en tres entregas progresivas, a las que se suma un simulacro de sustentación sin peso evaluativo.

\begin{itemize}
    \item Entrega 1 (formulación de la pregunta de negocio) vale el 0\% y es prerrequisito de las demás. Fecha límite: semana 6, del 31 de agosto al 5 de septiembre de 2026. Entrega a través de la plataforma Intu.
    \item Entrega 2 (base depurada y EDA) vale el 5\%. Sustentación: semana 9, del 21 al 26 de septiembre de 2026. Entrega a través de la plataforma Intu.
    \item Entrega 3 (sustentación final) vale el 20\%. Sustentación: semana 16, del 9 al 14 de noviembre de 2026. Entrega a través de la plataforma Intu.
    \item Simulacro de sustentación (semana 15, del 2 al 7 de noviembre de 2026): espacio de práctica previo a la sustentación final; no tiene peso evaluativo pero es obligatorio para recibir retroalimentación antes de la presentación definitiva.
\end{itemize}

\end{document}
